% ========== 2.3 USER REQUIREMENTS ==========
% Analysis of user needs and expectations
\clearpage
\section{2.3 User Requirements}
\addcontentsline{toc}{section}{2.3 User Requirements}
\label{sec:user-requirements}

\lettrine{U}{ser requirements} define what developers and other stakeholders need from the system to accomplish their goals effectively. These requirements emerge from analysis of current development workflows, identified problems with existing tools, and stakeholder expectations for improved productivity and collaboration.

Developers require seamless integration of AI capabilities into their existing workflows without forcing adoption of entirely new working methods. The system must support gradual adoption, allowing users to incorporate AI assistance at their own pace while maintaining access to familiar development tools and practices. Context preservation across different activities is essential, ensuring that the system maintains awareness of project state, developer intentions, and previous interactions without requiring repeated explanations.

Performance and responsiveness represent critical user requirements. Developers expect immediate feedback for all interactions, with response times comparable to traditional development tools. The system must provide clear status indicators for longer operations, remain responsive during background processing, and never leave users uncertain about system state. Resource consumption must remain within acceptable bounds, allowing developers to run other applications concurrently without performance degradation.

\textbf{Core User Requirements:}
\begin{compactlist}
    \item Natural interaction supporting both conversational and traditional editing workflows
    \item Transparent AI behavior with clear explanations of suggestions and decisions
    \item Complete human control with override capabilities at every decision point
    \item Adaptive learning from individual developer patterns and preferences
    \item Consistent quality in AI-generated outputs matching professional standards
\end{compactlist}

Users require transparency into AI decision-making processes and reasoning. When the system suggests code changes or implementation strategies, developers need to understand the rationale behind recommendations. This transparency enables informed decision-making and helps users learn from AI expertise rather than blindly accepting suggestions. Human override capabilities must be readily available, ensuring that developers maintain ultimate control over all development decisions and can correct system mistakes easily.
